% Work/MAT — přepracovaná hodina 15: Vlastnosti množinových operací.
% Pozor: nový hodinový rozvrh množin; původní hodina 15 se věnovala jen
% sjednocení a průniku. Po schválení bloku je třeba opravit ročníkový rozpis.
% Návaznost: 12 Co je množina, 13 Podmnožiny a operace,
% 14 Vennovy diagramy; příště operace s intervaly.
% Odborný zdroj: ucebnice/01_Zakladni_poznatky/Mnoziny/Množiny.tex,
% zejména Množinové operace, jejich vlastnosti a Vennovy diagramy.
% Metodická inspirace: Realisticky.cz, místní archiv:
% 010303_Mnozinove_operace/010303_Lekce.pdf — význam operací,
% 010304_Vennovy_diagramy/010304_Lekce.pdf — ověřování oblastí.
% Tyto externí podklady nejsou součástí dodaného artefaktu; příklady níže
% jsou nově sestavené a nevydáváme je za doslovný přepis jejich zadání.
% Všechny zákony pro doplněk se vztahují ke STEJNÉ základní množině U.
% Záměr: 45minutové jádro + rozšíření pro rychlé žáky a domácí práce.
\documentclass[11pt,a4paper]{article}
% Spolecny styl pro ucitelske pripravy MAT (Work/MAT)
% Verze 1.1, 2026-10-23
\NeedsTeXFormat{LaTeX2e}
\ProvidesPackage{priprava}[2026/10/23 v1.1 Sablona ucitelskych priprav MAT]

% Zakladni typografie a cestina
\RequirePackage[T1]{fontenc}
\RequirePackage[utf8]{inputenc}
\RequirePackage[czech]{babel}
\RequirePackage{lmodern}
\RequirePackage{microtype}
\RequirePackage{amsmath,amssymb}
\RequirePackage{xcolor}
\RequirePackage{enumitem}
\RequirePackage[most]{tcolorbox}
\RequirePackage{tabularx,booktabs,array}
\RequirePackage{fancyhdr}
\RequirePackage{needspace}
\RequirePackage{geometry}
\RequirePackage[hidelinks]{hyperref}

\geometry{a4paper,top=17mm,bottom=18mm,left=23mm,right=23mm}
\definecolor{notegray}{gray}{0.39}
\setlength{\parindent}{0pt}
\setlength{\parskip}{2.5pt plus 1pt minus 0.5pt}
\setlist[itemize]{topsep=2pt,itemsep=1pt,parsep=0pt,leftmargin=1.45em}
\setlist[enumerate]{topsep=2pt,itemsep=2pt,parsep=0pt,leftmargin=1.65em}
\setcounter{secnumdepth}{0}
\renewcommand{\arraystretch}{1.12}

% Zahlavi: vlevo rocnik | globalni tema, vpravo cislo: nazev hodiny
% Zapati zustava zcela prazdne.
\pagestyle{fancy}
\fancyhf{}
\renewcommand{\headrulewidth}{0.35pt}
\renewcommand{\footrulewidth}{0pt}
\setlength{\headheight}{14pt}
\newcommand{\nastavPripravu}[4]{%
  \fancyhead[L]{\small #1\enspace\textbar\enspace #2}%
  \fancyhead[R]{\small #3: #4}%
}

% Minimalisticke cernobile boxy. Kazda definice zvlast.
\tcbset{commonlesson/.style={enhanced,breakable=false,colback=white,colframe=black,
  colbacktitle=white,coltitle=black,boxrule=0.45pt,
  arc=0.8mm,outer arc=0.8mm,left=3mm,right=3mm,top=1.6mm,bottom=1.6mm,
  toptitle=1mm,bottomtitle=0.7mm,fonttitle=\bfseries,
  attach boxed title to top left={xshift=2.5mm,yshift=-1.15mm},
  boxed title style={colback=white,colframe=white,boxrule=0pt,left=1mm,right=1mm}}}
\newtcolorbox{zadaniBox}[1][]{commonlesson,title={Zadání},#1}
\newtcolorbox{prikladBox}[1][]{commonlesson,title={Příklad},#1}
\newtcolorbox{definiceBox}[1][]{commonlesson,title={Definice},#1}

% Spolecny box pro zadani a jeho reseni, bez mezery mezi castmi.
% Pouziti:
%   \begin{prikladReseniBox}
%     Zadani ulohy.
%     \tcblower
%     \begin{align*} ... \end{align*}
%   \end{prikladReseniBox}
% Spodni cast ma jemny sedy podklad a delici linku; lze v ni sazet
% vice radku matematiky. Obe casti maji jen jeden spolecny ramecek.
\tcbset{solutionlesson/.style={
  commonlesson,
  enhanced jigsaw,
  breakable,
  colbacklower=black!3,
  segmentation style={solid,draw=black!25,line width=0.35pt},
  fontlower=\small\color{notegray},
  before lower={\noindent\textbf{Řešení:}\enspace},
  leftlower=3mm,rightlower=3mm
}}
\newtcolorbox{prikladReseniBox}[1][]{solutionlesson,title={Příklad},#1}
\newtcolorbox{zadaniReseniBox}[1][]{solutionlesson,title={Zadání},#1}

% Pomucky pro autora pripravy: sedive poznamky; \reseni pro starsi pripravy
\newcommand{\poznamka}[1]{\par\begingroup\small\color{notegray}\itshape
  \textnormal{\textbf{Poznámka pro učitele.}}\enspace #1\par\endgroup}
\newcommand{\reseni}[1]{\par\begingroup\small\color{notegray}
  \textbf{Očekávané řešení:}\enspace #1\par\endgroup}
\newcommand{\blok}[1]{\needspace{5\baselineskip}\subsection*{#1}}
\newcommand{\vyrok}[1]{\uv{#1}}

\usepackage{tikz}
\usetikzlibrary{calc,patterns}
\newtcolorbox{vlastnostBox}[1][]{commonlesson,title={Vlastnost},#1}
\nastavPripravu{1. ročník}{Množiny}{15}{Vlastnosti množinových operací}

% Jednotné černobílé diagramy, stejné oblasti se značkují šedě.
\newcommand{\vennDve}[1]{%
\begin{tikzpicture}[scale=.68,baseline=(current bounding box.center)]
  \path[use as bounding box] (-2.72,-1.83) rectangle (2.72,1.79);
  #1
  \draw[thick] (-2.57,-1.63) rectangle (2.57,1.65);
  \draw[thick] (-.78,0) circle (1.25);
  \draw[thick] (.78,0) circle (1.25);
  \node at (-2.32,.15) {$A$}; \node at (2.32,.15) {$B$};
  \node[anchor=north east] at (2.49,1.56) {$U$};
\end{tikzpicture}}
\newcommand{\vennTri}[1]{%
\begin{tikzpicture}[scale=.68,baseline=(current bounding box.center)]
  \path[use as bounding box] (-2.72,-2.45) rectangle (2.72,2.05);
  #1
  \draw[thick] (-2.6,-2.35) rectangle (2.6,1.94);
  \draw[thick] (-.86,.37) circle (1.35);
  \draw[thick] (.86,.37) circle (1.35);
  \draw[thick] (0,-.76) circle (1.35);
  \node at (-2.35,.19) {$A$}; \node at (2.35,.19) {$B$};
  \node at (0,-2.24) {$C$};
  \node[anchor=north east] at (2.5,1.84) {$U$};
\end{tikzpicture}}
\newcommand{\sedaPrunik}{%
\begin{scope}\clip (-.78,0) circle (1.25);
\fill[black!26] (.78,0) circle (1.25);\end{scope}}
\newcommand{\sedaMimoSjednoceni}{%
\fill[black!23] (-2.57,-1.63) rectangle (2.57,1.65);
\fill[white] (-.78,0) circle (1.25);
\fill[white] (.78,0) circle (1.25);}
\newcommand{\sedaPrunikADvojice}{%
\begin{scope}\clip (-.86,.37) circle (1.35);
  \fill[black!24] (.86,.37) circle (1.35);
  \fill[black!24] (0,-.76) circle (1.35);
\end{scope}}
\newcommand{\sedaDvaPruniky}{%
\begin{scope}\clip (-.86,.37) circle (1.35);
 \clip (.86,.37) circle (1.35);
 \fill[black!24] (-2.6,-2.35) rectangle (2.6,1.94);\end{scope}
\begin{scope}\clip (-.86,.37) circle (1.35);
 \clip (0,-.76) circle (1.35);
 \fill[black!24] (-2.6,-2.35) rectangle (2.6,1.94);\end{scope}}

\begin{document}
\blok{Proč se učíme pravidla, když už umíme počítat?}
V hodině 13 jsme s množinami počítali, v hodině 14 jsme výsledky
vyznačovali do diagramu. Nyní zjišťujeme, které postupy lze
\textbf{bezpečně zaměnit, seskupit nebo zjednodušit}.
Stejně jako u čísel se vyplatí rozumět pravidlu, nejen si je pamatovat.

\poznamka{Jádro na 45 minut včetně organizace: 4 min motivace;
9 min komutativita, asociativita a P1; 12 min distributivita a P2;
10 min doplněk, De Morganovy zákony a P3; 7 min S1 a), b);
3 min shrnutí a zadání domácí práce. Důkaz libovolným prvkem vysvětli
u prvního distributivního zákona a u De Morganova zákona; další důkazy
jsou rozšíření. Neopisujte celou tabulku ze závěru.}

\blok{1. Na pořadí někdy záleží a někdy ne}
\begin{vlastnostBox}
\textbf{Komutativita} (záměna pořadí): pro libovolné množiny $A,B$ platí
\[
A\cup B=B\cup A,\qquad A\cap B=B\cap A.
\]
\textbf{Asociativita} (změna seskupení): pro libovolné $A,B,C$ platí
\[
(A\cup B)\cup C=A\cup(B\cup C),\quad
(A\cap B)\cap C=A\cap(B\cap C).
\]
\end{vlastnostBox}

\begin{prikladReseniBox}[title={P1 — jádro: co lze prohodit?}]
Nechť $A=\{1,2,4\}$ a $B=\{2,3,4\}$.
Vypočti $A\cup B$, $B\cup A$, $A\cap B$, $B\cap A$
a oba rozdíly $A\setminus B$, $B\setminus A$.
\tcblower
\begin{align*}
A\cup B&=B\cup A=\{1,2,3,4\},\\[-2pt]
A\cap B&=B\cap A=\{2,4\},\\[-2pt]
A\setminus B&=\{1\},\qquad B\setminus A=\{3\}.
\end{align*}
Sjednocení a průnik jsou komutativní, ale \textbf{rozdíl obecně není}.
Tento výpočet ukazuje pravidla na jedné dvojici; sám o sobě nedokazuje,
že platí pro všechny množiny.
\end{prikladReseniBox}

\poznamka{Přirovnání k logice: u výroku \uv{patří do $A$ nebo do $B$}
nezáleží na pořadí podmínek, podobně u \uv{a zároveň}. Naopak
\uv{je v $A$, ale není v $B$} se po prohození může změnit.
Asociativitu stačí v jádru vysvětlit přeskupením tří podmínek;
neslučuj ji s chybným přesunem závorek mezi různými operacemi.}

\blok{2. Distributivita: jedna podmínka pro dvě možnosti}
\begin{vlastnostBox}
\textbf{Distributivita} propojuje sjednocení a průnik:
\[
\begin{aligned}
A\cap(B\cup C)&=(A\cap B)\cup(A\cap C),\\
A\cup(B\cap C)&=(A\cup B)\cap(A\cup C).
\end{aligned}
\]
\end{vlastnostBox}

\begin{prikladReseniBox}[title={P2 — jádro: obě strany vycházejí stejně}]
Pro $A=\{1,2,4,6\}$, $B=\{2,3,6\}$ a $C=\{1,3,6\}$
nezávisle spočti obě strany prvního distributivního zákona.
\tcblower
Levá strana: $B\cup C=\{1,2,3,6\}$, takže
$A\cap(B\cup C)=\{1,2,6\}$.

Pravá strana: $A\cap B=\{2,6\}$ a $A\cap C=\{1,6\}$,
proto $(A\cap B)\cup(A\cap C)=\{1,2,6\}$.
\end{prikladReseniBox}

\blok{Důkaz: proč distributivita platí pro libovolné množiny?}
U příkladu jsme ověřili jedinou trojici množin.
Teď si vyberme \textbf{libovolný prvek} $x$:
\[
\begin{aligned}
x\in A\cap(B\cup C)
&\Longleftrightarrow x\in A\ \land\ (x\in B\ \lor\ x\in C)\\
&\Longleftrightarrow (x\in A\ \land\ x\in B)\ \lor\
                  (x\in A\ \land\ x\in C)\\
&\Longleftrightarrow x\in(A\cap B)\cup(A\cap C).
\end{aligned}
\]
Proto mají obě množiny právě stejné prvky — jde o \textbf{důkaz}.

\textbf{Stejná oblast, dva postupy:}
\begin{center}
\begin{tabularx}{\linewidth}{@{}XX@{}}
\centering\vennTri{\sedaPrunikADvojice} &
\centering\vennTri{\sedaDvaPruniky}\tabularnewline[-.2em]
\centering\small $A\cap(B\cup C)$ &
\centering\small $(A\cap B)\cup(A\cap C)$\tabularnewline
\end{tabularx}
\end{center}

\poznamka{Ať žák nejprve slovně popíše levé vybarvení: \uv{leží v $A$
a současně alespoň v jednom z $B,C$}. Teprve pak porovnej se dvěma
překryvy napravo. Druhý distributivní zákon pouze ukaž; jeho důkaz
je v rezervě.}\par

\blok{3. Doplněk a De Morganovy zákony}
Od této chvíle jsou všechny množiny podmnožinami \textbf{stejného}
$U$ a $A'=U\setminus A$, $B'=U\setminus B$.

\begin{vlastnostBox}
\textbf{Doplněk:} $A\cup A'=U$, $A\cap A'=\varnothing$, $(A')'=A$.

\textbf{De Morganovy zákony} (\uv{negace obrací operaci}):
\[
(A\cup B)'=A'\cap B',\qquad
(A\cap B)'=A'\cup B'.
\]
\end{vlastnostBox}

\textbf{Dva zápisy stejné šedé oblasti:}
\begin{center}
\begin{tabularx}{\linewidth}{@{}XX@{}}
\centering\vennDve{\sedaMimoSjednoceni} &
\centering\vennDve{\sedaMimoSjednoceni}\tabularnewline[-.3em]
\centering\small $(A\cup B)'$ & \centering\small $A'\cap B'$\tabularnewline
\end{tabularx}
\end{center}

\begin{prikladReseniBox}[title={P3 — jádro: doplněk sjednocení}]
Nechť $U=\{1,2,\ldots,8\}$, $A=\{1,2,4,7\}$,
$B=\{2,3,7,8\}$. Ověř oba De Morganovy zákony výčtem.
\tcblower
$A'=\{3,5,6,8\}$, $B'=\{1,4,5,6\}$.
\begin{align*}
A\cup B&=\{1,2,3,4,7,8\},\\
(A\cup B)'&=\{5,6\}=A'\cap B'.
\end{align*}
Dále $A\cap B=\{2,7\}$, a tedy
\[
(A\cap B)'=\{1,3,4,5,6,8\}=A'\cup B'.
\]
Obecný důvod prvního zákona: aby prvek \textbf{nebyl ani v jedné}
z množin $A,B$, musí být \textbf{mimo $A$ a zároveň mimo $B$}.
Proto patří právě do $A'\cap B'$.
\end{prikladReseniBox}

\poznamka{Vezmi známé logické pravidlo \uv{není pravda, že $P$ nebo $Q$}
znamená \uv{neplatí $P$ a neplatí $Q$}. De Morganovy zákony tedy
navazují na výrokovou logiku. Upozorni na častý omyl:
$(A\cup B)'\ne A'\cup B'$ obecně.}

\blok{4. Samostatné procvičení — jádro}
\begin{zadaniReseniBox}[title={S1 — rozhodni a vypočti}]
\begin{enumerate}[label=\alph*)]
\item Rozhodni, která tvrzení platí \emph{pro všechny} množiny
$A,B\subseteq U$. Neplatné tvrzení oprav:
\[
A\cap U=A;\qquad A\cup\varnothing=\varnothing;\qquad
(A\cup B)'=A'\cup B';\qquad (A')'=A.
\]
\item Pro $U=\{0,1,\ldots,6\}$, $A=\{0,1,2,5\}$,
$B=\{1,3,5\}$, $C=\{0,3,5\}$ ověř výčtem
$A\cap(B\cup C)=(A\cap B)\cup(A\cap C)$.
\item (Volitelně) Slovně vysvětli, proč je $A\cap\varnothing=\varnothing$.
\end{enumerate}
\tcblower
\begin{enumerate}[label=\alph*)]
\item První ano. Druhé ne, správně $A\cup\varnothing=A$.
Třetí obecně ne, správně $(A\cup B)'=A'\cap B'$.
Čtvrté ano, pokud oba doplňky bereme ve stejném $U$.
\item $B\cup C=\{0,1,3,5\}$, takže levá strana je $\{0,1,5\}$.
$A\cap B=\{1,5\}$ a $A\cap C=\{0,5\}$,
jejich sjednocení je také $\{0,1,5\}$.
\item Žádný prvek nemůže současně patřit do $A$ a do prázdné množiny.
\end{enumerate}
\end{zadaniReseniBox}

\poznamka{V povinném jádru stačí a), b); c) je rychlá kontrolní otázka.
Nech žáka říci, zda zákon jen ověřil na datech, nebo zdůvodnil obecně.
Neplatný obecný zákon vyvrátí i jediný protipříklad.}

\blok{5. Rozšíření — užitečné vztahy a důkazy}
\begin{vlastnostBox}
\textbf{Idempotence:} $A\cup A=A$, $A\cap A=A$.
\textbf{Absorpce:} $A\cup(A\cap B)=A$, $A\cap(A\cup B)=A$.
\textbf{Prázdná a základní množina:}
\[
A\cup\varnothing=A,\quad A\cap\varnothing=\varnothing,\quad
A\cup U=U,\quad A\cap U=A.
\]
\end{vlastnostBox}

\begin{prikladReseniBox}[title={R1 — dokážeme absorpci, ne jen vyzkoušíme}]
Vysvětli, proč $A\cup(A\cap B)=A$ pro \emph{libovolné} množiny.
Ukaž jeden protipříklad k tvrzení
$A\setminus(B\cup C)=(A\setminus B)\cup(A\setminus C)$.
\tcblower
Každý prvek $A\cap B$ již patří do $A$. Jeho přidáním do $A$
nevznikne žádný nový prvek; sjednocení proto zůstane $A$.

Protipříklad: $A=\{1\}$, $B=\{1\}$, $C=\varnothing$.
Vlevo $A\setminus(B\cup C)=\varnothing$,
vpravo $(A\setminus B)\cup(A\setminus C)=\{1\}$.
Správný vztah má vpravo \textbf{průnik}:
$A\setminus(B\cup C)=(A\setminus B)\cap(A\setminus C)$.
\end{prikladReseniBox}

\begin{prikladReseniBox}[title={R2 — druhý De Morganův zákon a distributivita}]
Pomocí podmínek příslušnosti libovolného prvku $x$ zdůvodni
$(A\cap B)'=A'\cup B'$. Potom analogicky zdůvodni
$A\cup(B\cap C)=(A\cup B)\cap(A\cup C)$.
\tcblower
$x\notin A\cap B$ právě tehdy, když $x\notin A$ \textbf{nebo}
$x\notin B$; to je příslušnost k $A'\cup B'$.

U druhé rovnosti: pokud $x\in A$, patří do obou stran.
Jinak je na levé straně právě tehdy, když $x\in B\cap C$;
totéž pak vyžadují obě části pravého průniku.
\end{prikladReseniBox}

\blok{6. Domácí procvičení}
\begin{zadaniReseniBox}[title={D1 — pravidlo, výpočet a zdůvodnění}]
Pro a) až c) je $U=\{1,2,\ldots,10\}$,
$A=\{2,4,6,8,10\}$, $B=\{1,2,3,4\}$, $C=\{3,6,9\}$.
\begin{enumerate}[label=\alph*)]
\item Vypočti obě strany $A\cap(B\cup C)=(A\cap B)\cup(A\cap C)$.
\item Vypočti $(A\cup B)'$ a ověř výsledek jako $A'\cap B'$.
\item Vypočti $(A\cap B)'$ a ověř výsledek jako $A'\cup B'$.
\item Rozhodni, zda platí $A\setminus B=B\setminus A$ pro každé $A,B$.
Uveď protipříklad.
\item Zjednoduš $M\cap(M\cup N)$ a výsledek zdůvodni slovně.
\end{enumerate}
\tcblower
\begin{enumerate}[label=\alph*)]
\item $B\cup C=\{1,2,3,4,6,9\}$ a
$A\cap(B\cup C)=\{2,4,6\}$.
$A\cap B=\{2,4\}$, $A\cap C=\{6\}$,
jejich sjednocení je $\{2,4,6\}$.
\item $A\cup B=\{1,2,3,4,6,8,10\}$, proto
$(A\cup B)'=\{5,7,9\}$.
$A'=\{1,3,5,7,9\}$, $B'=\{5,6,7,8,9,10\}$,
$A'\cap B'=\{5,7,9\}$.
\item $A\cap B=\{2,4\}$ a
$(A\cap B)'=\{1,3,5,6,7,8,9,10\}$.
Stejnou množinu dostaneme jako $A'\cup B'$.
\item Ne, například pro $A=\{1,2\}$, $B=\{2\}$ je
$A\setminus B=\{1\}$, zatímco $B\setminus A=\varnothing$.
\item $M$: protože $M\subseteq M\cup N$, průnik s větší množinou
neodebere žádný prvek $M$.
\end{enumerate}
\end{zadaniReseniBox}

\blok{Rychlý přehled pro učitele (volitelná závěrečná tabulka)}
\small
\begin{tabularx}{\linewidth}{@{}>{\bfseries}p{25mm}X@{}}
\toprule
Název & Vztahy, které si mají žáci odnést \\ \midrule
Záměna & $A\cup B=B\cup A$; $A\cap B=B\cap A$ \\
Seskupení & $(A\cup B)\cup C=A\cup(B\cup C)$ a obdobně pro $\cap$ \\
Roznásobení & $A\cap(B\cup C)=(A\cap B)\cup(A\cap C)$;\newline
$A\cup(B\cap C)=(A\cup B)\cap(A\cup C)$ \\
De Morgan & $(A\cup B)'=A'\cap B'$; $(A\cap B)'=A'\cup B'$ \\
Doplněk & $A\cup A'=U$, $A\cap A'=\varnothing$, $(A')'=A$ \\
\bottomrule
\end{tabularx}
\normalsize

\poznamka{Závěrečná otázka: Proč platnost dvou stran zákona pro jedny
konkrétní množiny ještě není důkaz? Protože obecná rovnost tvrdí
shodu pro \emph{všechny} přípustné množiny. K vyvrácení stačí jeden
protipříklad. Na další hodině tyto operace použijeme na intervaly.}

\blok{K rychlému vyznačení při promítání (volitelné)}
Vyšrafuj vlevo $(A\cap B)'$ a vpravo $A\cap(B\cup C)$.
Pak řekni, které jiné zápisy popisují stejné oblasti.
\begin{center}
\begin{tabularx}{\linewidth}{@{}XX@{}}
\centering\vennDve{} & \centering\vennTri{}\tabularnewline
\end{tabularx}
\end{center}

% -----------------------------------------------------------------------------
% VOLITELNE ROZSIRENI: mimo 45minutove jadro. Vychazi z oddilu Mohutnost mnoziny
% a Nekonecne mnoziny v projektove ucebnici; matematicka zpracovani jsou vlastni.
% -----------------------------------------------------------------------------
\clearpage
\blok{7. BONUS: Mohutnost množiny -- kolik má prvků?}
\poznamka{\textbf{Rozšíření mimo povinných 45 minut.} První dva příklady lze
použít jako krátkou pětiminutovou vsuvku, celý výklad nekonečen spíše jako
matematickou zajímavost, dobrovolný seminář nebo nadstavbu pro zvídavé žáky.
Není určen k opisování ani k běžnému zkoušení.}

\begin{definiceBox}
\textbf{Mohutnost (kardinalita) konečné množiny} $A$ je počet jejích
\emph{různých} prvků; značíme ji $|A|$. Zejména $|\varnothing|=0$.
\end{definiceBox}

\begin{prikladReseniBox}[title={B1 -- počet prvků není délka zápisu}]
Urči mohutnosti množin
$A=\{1,2,2,3,1\}$, $B=\{0,2,4,6\}$,
$C=\varnothing$ a $D=\{\varnothing\}$.
\tcblower
$A=\{1,2,3\}$, tedy $|A|=3$; $|B|=4$ a $|C|=0$.
Množina $D$ obsahuje \emph{jeden prvek}, kterým je prázdná množina,
proto $|D|=1$. Opakované zapsání prvku jeho počet nezvyšuje.
\end{prikladReseniBox}

\blok{8. Jak porovnat dvě nekonečné množiny?}
U konečných množin prvky spočítáme. U nekonečných to nejde dokončit.
Přesto můžeme porovnat jejich \emph{velikost} jinak: každému prvku jedné
množiny přiřadíme právě jeden prvek druhé a na žádný nezapomeneme.

\begin{definiceBox}
Množiny $A$ a $B$ mají \textbf{stejnou mohutnost}, jestliže mezi nimi
existuje \textbf{bijekce}: vzájemně jednoznačné přiřazení, v němž má každý
prvek $A$ právě jednoho partnera v $B$ a každý prvek $B$ právě jednoho
partnera v $A$.
\end{definiceBox}

\textbf{Rozcvička:} Přirozená a sudá kladná čísla můžeme spárovat jako
$1\leftrightarrow2$, $2\leftrightarrow4$, $3\leftrightarrow6$, $\ldots$,
obecně $n\leftrightarrow2n$. Žádné sudé kladné číslo nezůstane bez partnera!

\begin{prikladReseniBox}[title={B2 -- překvapení: $|\mathbb N|=|\mathbb Z|$}]
Má množina všech celých čísel více prvků než množina přirozených čísel?
Používáme $\mathbb N=\{1,2,3,\ldots\}$.
\tcblower
Celá čísla můžeme seřadit takto:
\[
\boxed{0,\ 1,\ -1,\ 2,\ -2,\ 3,\ -3,\ 4,\ -4,\ \ldots}
\]
A pak každému přirozenému číslu přiřadit jeho položku v tomto seznamu:
\begin{center}
\begin{tikzpicture}[x=1.56cm,y=.82cm,>=stealth,
  every node/.style={font=\normalsize}]
  \node[anchor=east] at (-.4,0) {$\mathbb N:$};
  \node[anchor=east] at (-.4,-1.7) {$\mathbb Z:$};
  \foreach \i/\a/\b in {0/1/0,1/2/1,2/3/-1,3/4/2,4/5/-2,5/6/3,6/7/-3}{
    \node (n\i) at (\i,0) {$\a$};
    \node (z\i) at (\i,-1.7) {$\b$};
    \draw[->,thin] (n\i.south) -- (z\i.north);
  }
  \node at (7.15,-.85) {$\cdots$};
\end{tikzpicture}
\end{center}
\textbf{Proč se to nikdy nezasekne?} Po nule postupně střídáme
kladná a záporná celá čísla se stále větší absolutní hodnotou.
Každé celé číslo se jednou objeví a žádné se neopakuje.

Dokonce máme přesný předpis (pro $k=0,1,2,\ldots$):
\[
f(2k+1)=-k,\qquad f(2k)=k\quad (k\ge1).
\]
Každé přirozené číslo tedy označí právě jedno celé číslo a
každé celé číslo získá právě jedno pořadové číslo.
Proto \textbf{$|\mathbb N|=|\mathbb Z|$}, přestože $\mathbb N\subsetneq\mathbb Z$.
\end{prikladReseniBox}

\poznamka{U nekonečných množin může mít vlastní podmnožina stejnou
mohutnost. Konečný obrázek se šipkami je ilustrace; rozhoduje
bijekce určená obecným pravidlem.}

\blok{9. BONUS: Jsou všechna nekonečna stejně velká?}
\begin{definiceBox}
\textbf{Spočetně nekonečná množina} má stejnou mohutnost jako
$\mathbb N$. Její mohutnost zapisujeme $\aleph_0$ (čteme
\uv{alef nula}). \textbf{Nespočetná množina} nemůže být takto očíslována
přirozenými čísly.
\end{definiceBox}

\begin{center}
\begin{tabularx}{\linewidth}{@{}>{\raggedright\arraybackslash}p{35mm}X@{}}
\toprule
Množina & Mohutnost a proč \\ \midrule
$\mathbb N$ & $\aleph_0$: základní nekonečný seznam $1,2,3,\ldots$ \\
$\mathbb Z$ & $\aleph_0$: explicitní párování v příkladu B2 \\
$\mathbb Q$ & $\aleph_0$: zlomky $p/q$ lze systematicky procházet po skupinách
(vynecháváme opakované zápisy téhož čísla) \\
$\mathbb R$ & $\mathfrak c=2^{\aleph_0}$: \emph{kontinuum}, nespočetná mohutnost \\
$(0;1)$, $\mathbb R\setminus\mathbb Q$ & rovněž $\mathfrak c$, ačkoli jde o vlastní podmnožiny $\mathbb R$ \\
\bottomrule
\end{tabularx}
\end{center}

\begin{prikladReseniBox}[title={B3 -- proč se mezi zlomky dá postupovat systematicky?}]
Racionální čísla vznikají ze zlomků $p/q$, kde $p\in\mathbb Z$
a $q\in\mathbb N$. Proč je můžeme postupně všechny vyjmenovat,
i když jich je mezi dvěma různými čísly nekonečně mnoho?
\tcblower
Procházíme dvojice $(p,q)$ podle rostoucí hodnoty $|p|+q$.
Pro každou pevnou hodnotu součtu existuje jen konečně mnoho takových
dvojic. Proto se po konečně mnoha krocích dostaneme ke každému
konkrétnímu zlomku. Pokud už jsme stejnou hodnotu získali jiným
zlomkem, podruhé ji nezařadíme. Tak lze $\mathbb Q$ očíslovat a
$|\mathbb Q|=\aleph_0$.
\end{prikladReseniBox}

\begin{prikladReseniBox}[title={B4 -- Cantorův diagonální argument, intuitivně}]
Představme si, že někdo tvrdí, že sepsal \emph{všechna} reálná čísla
z intervalu $(0;1)$ do nekonečného seznamu $r_1,r_2,r_3,\ldots$.
Jak ukázat, že v něm nějaké číslo určitě chybí?
\tcblower
U každého čísla zapiš desetinný rozvoj (zvolíme zápis, který
nekončí nekonečnou řadou devítek). Vezmi první desetinnou číslici
prvního čísla, druhou druhého, třetí třetího atd.

Sestroj nové číslo $s=0,s_1s_2s_3\ldots$ tak, že $s_n=2$, pokud
má $r_n$ na $n$-tém desetinném místě číslici $1$;
ve všech ostatních případech polož $s_n=1$.

Nové číslo se \textbf{od $r_n$ liší na $n$-tém místě}, a tedy
\emph{není rovno žádnému} číslu ze seznamu. Přitom leží v $(0;1)$.
Údajně úplný seznam tedy nemohl existovat.
Proto $|\mathbb R|>\aleph_0$.
\end{prikladReseniBox}

\poznamka{Pro zvědavé: $\aleph_1$ označuje \emph{nejmenší nespočetnou
kardinální mohutnost}; platí $\aleph_0<\aleph_1\le\mathfrak c$.
\emph{Hypotéza kontinua} tvrdí $\mathfrak c=\aleph_1$.
V běžné axiomatické teorii množin ZFC ji nelze dokázat ani vyvrátit
(pokud je ZFC bezesporná). Proto \textbf{nepiš automaticky}
$|\mathbb R|=\aleph_1$. Tato poslední poznámka je pouze pro zájemce.}

\textbf{Dvě otázky pro zájemce.} Proč má $\{\varnothing\}$ jeden prvek,
i když je $\varnothing$ prázdná? A které přirozené číslo při párování
z B2 odpovídá celému číslu $-4$?
\poznamka{Řešení: jediným prvkem $\{\varnothing\}$ je právě prázdná
množina, takže $|\{\varnothing\}|=1$. Číslo $-4$ je devátou položkou
seznamu ($f(9)=-4$). Mohutnost neměníme pouhým přeuspořádáním prvků.
Pro hodinu 16 není tato nadstavba předpokladem.}
\end{document}
